\chapter{Metodologia}
\label{cap:metodologia}

\section{Classificação da Pesquisa}
\label{sec:classificacao-da-pesquisa}

natureza → aplicada
abordagem → quantitativa
objetivo → exploratória / descritiva
método → experimental

\section{Estratégia de Pesquisa}
\label{sec:estrategia-de-pesquisa}

estudo comparativo
análise entre duas arquiteturas
uso de sistemas reais
unidade de análise (classes, pacotes, etc.)

\section{Definição das Métricas}
\label{sec:definicao-das-metricas}

quais métricas vai usar (CBO, LCOM, etc.)
o que cada uma mede
por que são relevantes

\section{Ferramentas Utilizadas}
\label{sec:ferramentas-utilizadas}

ferramentas de extração (ex: CK, SonarQube)
ambiente (Java, etc.)
como elas são usadas (sem entrar em passo a passo)

\section{Procedimento Experimental}
\label{sec:procedimento-experimental}

Aqui é o “passo a passo”:

seleção dos sistemas
execução das ferramentas
coleta das métricas
organização dos dados

\section{Tratamento e Análise dos Dados}
\label{sec:tratamento-e-analise-dos-dados}

como vai comparar os sistemas
médias, distribuições, etc.
normalização (se necessário)
critérios de comparação

\section{Ameaças à Validade}
\label{sec:ameacas-a-validez}

Você discute possíveis limitações:

diferença de tamanho dos sistemas
viés dos projetos (foram feitos por você)
generalização dos resultados
limitações das métricas